% SGA_NATIVE_CONTINUATION_389_390
\providecommand{\Sym}{\operatorname{Sym}}
\providecommand{\GL}{\operatorname{GL}}
\providecommand{\Grass}{\operatorname{Grass}}
\providecommand{\gr}{\operatorname{gr}}
\providecommand{\ctop}{c_{\mathrm{top}}}
\providecommand{\Rep}{\operatorname{Rep}}
\providecommand{\RZ}{R_{\mathbb Z}}
\providecommand{\Dflag}{D}
\providecommand{\Gmn}{G}
complete the induction, noting that, if @@M0@@, and if @@M1@@ for @@M2@@, we have:
@@M3@@
whence:
@@M4@@
whence @@M5@@ by what precedes.

\sgaStableTarget{sga6.e06.sec.3}\section*{3. Consequences of the structure theorem for the @@M6@@ of a projective bundle}

\sgaStableTarget{sga6.e06.prop.3-1}\textbf{Proposition 3.1.} \uline{Let} @@M7@@ \uline{be a quasi}-\uline{compact scheme}, @@M8@@ \uline{locally free} @@M9@@\uline{-Modules of finite type}, @@M10@@ \uline{the flag scheme of} @@M11@@, @@M12@@.
\uline{Then the canonical homomorphism}:
@@M13@@
\uline{is injective} (@@M14@@ \uline{being defined}, \uline{either from locally free} @@M15@@\uline{-Modules of finite type}, \uline{or from perfect complexes}).

We check immediately, using the universal property of flag schemes, that @@M16@@ can be constructed as follows: we consider the sequence of @@M17@@-schemes @@M18@@, defined inductively by taking @@M19@@ to be the flag scheme of the inverse image of @@M20@@ over @@M21@@; then @@M22@@. We then see at once that, to prove the proposition, it suffices to prove it in the case @@M23@@, the general case following by induction. We can therefore assume that @@M24@@ is the flag scheme of a locally free @@M25@@-Module of finite type @@M26@@.

We know that @@M27@@ can then be constructed as follows: we put @@M28@@, @@M29@@, and define inductively schemes @@M30@@ and
